% ==================================================================
% house-style.tex — Safe Notation & Helpers
% ==================================================================

\makeatletter % <--- РАЗРЕШАЕМ @

% ----- Small helpers (Delimiters) -----
\providecommand{\abs}[1]{\left\lvert #1\right\rvert}
\providecommand{\norm}[1]{\left\lVert #1\right\rVert}
\providecommand{\ip}[2]{\left\langle #1,\,#2\right\rangle}
\providecommand{\set}[1]{\left\{\,#1\,\right\}}

% ----- Symbols -----
\providecommand{\dd}{\mathrm{d}}
\providecommand{\eps}{\varepsilon}
\providecommand{\OO}{\mathcal{O}}
\providecommand{\oo}{o}

% ----- Number sets -----
\providecommand{\R}{\mathbb{R}}
\providecommand{\C}{\mathbb{C}}
\providecommand{\Z}{\mathbb{Z}}
\providecommand{\N}{\mathbb{N}}
\providecommand{\Hh}{\mathbb{H}}

% ----- Operators (Guarded) -----
% Исправленный хелпер: правильно создает команду из текста
\providecommand{\SafeDeclareOp}[2]{%
  \@ifundefined{#1}{%
    \expandafter\DeclareMathOperator\csname #1\endcsname{#2}%
  }{}%
}

\SafeDeclareOp{inj}{inj}
\SafeDeclareOp{diam}{diam}

% Повторяем квантование на случай, если operators.tex не загружен
\SafeDeclareOp{Op}{Op}
\SafeDeclareOp{Opw}{Op_h^{\mathrm{w}}}
\SafeDeclareOp{Opl}{Op_h^{\mathrm{left}}}
\SafeDeclareOp{Opr}{Op_h^{\mathrm{right}}}

% ----- TODO helper -----
\@ifundefined{TODO}{
  \newcommand{\TODO}[1]{\textcolor{red!70!black}{\textbf{TODO:}~#1}}
}{}

\makeatother % <--- ЗАКРЫВАЕМ @
